%% =====================================================================
%%  T-07-03  The Three-Second Read
%%  -------------------------------------------------------------------
%%  One idea per file. Core is mandatory. Slots in canonical order.
%%  Max 700 words. No layout commands. No filler. New source material
%%  for this entry goes in sources/<BOOK>/T-07-03.tex -- never by
%%  rewriting this file (docs/RULES.md R0.1).
%% =====================================================================
%% @kind: topic
%% @id: T-07-03
%% @title: The Three-Second Read
%% @subtitle: The involuntary response precedes the managed answer
%% @chapter: c07
%% @order: 30
%% @status: stable
%% @sources: B1
%% @updated: 2026-09-19

\topic{T-07-03}{The Three-Second Read}{The involuntary response precedes the managed answer}

\begin{core}
Between a question and its answer there is a short window in which the face has not yet been arranged. What appears in that window is closer to the person's actual position than anything they subsequently say.
\end{core}
\begin{mechanism}
Answering takes work: the content must be selected, the framing chosen, the delivery managed. That work takes a second or two, and before it completes, the immediate affective response to the question is already visible --- surprise, recognition, fear, relief, anger. The response is to the question, not to the answer being constructed, which is what makes it informative. Reading it requires that the question be asked directly and suddenly enough that no prepared answer exists, and that the observer be looking at the right place at the right moment rather than listening for the words.
\end{mechanism}
\begin{tells}
  \item A micro-expression of recognition when a name or subject is introduced.
  \item A beat of silence before an answer that should have been immediate.
  \item The eyes and upper face move before the mouth does.
  \item Relief rather than engagement when the subject is dropped.
  \item An answer that does not match the expression that preceded it.
\end{tells}
\begin{moves}
  \item Introduce the subject indirectly first, then ask the direct question once, cleanly.
  \item Look at the eye region and upper cheekbones, not the mouth.
  \item Read the first three seconds and then stop reading. Later expressions are managed.
  \item Compare the expression against the answer. Mismatch is the finding.
  \item Do not follow up immediately. The value of the read is lost if the counterpart knows it happened.
\end{moves}
\begin{counters}
  \item Assume your face answers before you do, particularly on subjects that matter to you.
  \item Introduce a deliberate pause before answering anything unexpected; a uniform delay is uninformative in a way that a natural reaction is not.
  \item Answer the first version of the question and let the follow-up arrive later, when you have had time.
  \item Where you are watched closely, control the setting rather than the face: move, eat, face a light, involve other people.
  \item Do not over-trust your own reads. You are much worse at this than you believe, and confident misreading is worse than no reading.
\end{counters}
\begin{cost}
This is the least reliable entry in the chapter and the one most prone to confident error. Expression reading has a poor evidence base in general; cultural, individual and situational variation is large; and the observer's expectations generate the readings they expect to see. The source presents it as dependable, which is not supported. Its usable form is narrow: a mismatch between a clear expression and a clear answer, on a subject you already know matters, from someone you observe often. Outside that it should not change a decision on its own.
\end{cost}

\sources{B1}
\seesources
\seealso{T-06-02, T-07-01, T-07-02}
